\documentclass{article}
\usepackage{amsmath,physics,ctex}
\usepackage{bm}
\usepackage[top=1.7cm, bottom=1.5cm, left=1.8cm, right=1.8cm]{geometry}
\begin{document}
	\begin{center}
		\large{\textbf{题目}}
	\end{center}
	\medskip
	证明如下三个有关球谐函数$Y_{l}^{m}(\theta, \phi)$的等式，其中$\delta_{i,\, j}$是KroneckerDelta符号，上加横线表示复共轭。
	\begin{align*}
		&\int_0^\pi \!\! \int_0^{2\pi}
		\overline{Y_{L}^{M}(\theta, \phi)}\,
		e^{i\phi}
		\left(
		\pdv{Y_{l}^{m}(\theta, \phi)}{\theta} \sin\theta \cos\theta +
		i \pdv{Y_{l}^{m}(\theta, \phi)}{\phi}
		\right)
		 d\phi\, d\theta \\
		&= \delta_{M,\, m+1}
		\left[
		\delta_{L,\, l+1}\, l\, \sqrt{\frac{(l + m + 1)(l + m + 2)}{(2l + 1)(2l + 3)}}
		+ \delta_{L,\, l-1}\, (l + 1)\, \sqrt{\frac{(l - m)(l - m - 1)}{4l^2 - 1}}
		\right];
		\\[1em]
		&\int_0^\pi \!\! \int_0^{2\pi}
		\overline{Y_{L}^{M}(\theta, \phi)}\,
		e^{-i\phi}
		\left(
		\pdv{Y_{l}^{m}(\theta, \phi)}{\theta} \sin\theta \cos\theta -
		i \pdv{Y_{l}^{m}(\theta, \phi)}{\phi}
		\right)
		 d\phi\, d\theta \\
		&= -\delta_{M,\, m-1}
		\left[
		\delta_{L,\, l+1}\, l\, \sqrt{\frac{(l - m + 1)(l - m + 2)}{(2l + 1)(2l + 3)}}
		+ \delta_{L,\, l-1}\, (l + 1)\, \sqrt{\frac{(l + m)(l + m - 1)}{4l^2 - 1}}
		\right];
		\\[1em]
		&\int_0^\pi \!\! \int_0^{2\pi}
		\overline{Y_{L}^{M}(\theta, \phi)}\,
		\pdv{Y_{l}^{m}(\theta, \phi)}{\theta}
		\sin^2\theta\, d\phi\, d\theta \\
		&= \delta_{M,\, m}
		\left[
		\delta_{L,\, l+1}\, l\, \sqrt{\frac{(l + 1)^2 - m^2}{(2l + 1)(2l + 3)}}
		- \delta_{L,\, l-1}\, (l + 1)\, \sqrt{\frac{l^2 - m^2}{4l^2 - 1}}
		\right];
	\end{align*}
	
	并依此求解如下球坐标下的积分在直角坐标系的分量形式。
	\[
	\int_0^\pi \!\! \int_0^{2\pi}
	\overline{Y_{L}^{M}(\theta, \phi)}\,
	r\,(\nabla{Y_{l}^{m}(\theta, \phi)})\sin\theta\, d\phi\, d\theta \\
	\]
	\bigskip
	\begin{center}
		\large{\textbf{解答}}
	\end{center}
	\medskip
	根据连带勒让德函数递推关系：
	\begin{align*}
		x P_l^m(x) &= \frac{l - m + 1}{2l + 1} P_{l + 1}^m(x) + \frac{l + m}{2l + 1} P_{l - 1}^m(x), \\
		\sqrt{1 - x^2}P_l^m(x) &= \frac{1}{2l + 1}\left(-P_{l + 1}^{m + 1}(x) + P_{l - 1}^{m + 1}(x)\right), \\
		\sqrt{1 - x^2}P_l^m(x) &= \frac{1}{2l + 1}\left((l - m + 1)(l - m + 2)P_{l + 1}^{m + 1}(x) - (l + m - 1)(l + m)P_{l - 1}^{m + 1}(x)\right).
	\end{align*}
	
	易得：
	\begin{align*}
		\int_{-1}^{1} P_L^m(x) P_l^m(x) x \,dx 
		&\quad= \frac{2(l+m+1)!}{(2l+1)(2l+3)(l-m)!}\delta_{L,l+1} + \frac{2(l+m)!}{(4l^2-1)(l-m-1)!}\delta_{L,l-1}, \\
		m\int_{-1}^{1} \sqrt{1 - x^2} P_l^m(x) P_L^{m+1}(x) \,dx 
		&\quad= -\frac{2m(l+m+2)!}{(2l+1)(2l+3)(l-m)!}\delta_{L,l+1} + \frac{2m(l+m)!}{(4l^2-1)(l-m-2)!}\delta_{L,l-1}, \\
		m\int_{-1}^{1} \sqrt{1 - x^2} P_l^m(x) P_L^{m-1}(x) \,dx 
		&\quad= \frac{2m(l+m)!}{(2l+1)(2l+3)(l-m)!}\delta_{L,l+1} - \frac{2m(l+m)!}{(4l^2-1)(l-m)!}\delta_{L,l-1}.
	\end{align*}
	
	根据连带勒让德函数导数的递推关系：
	\begin{align*}
		&(1 - x^2)\frac{dP_l^m(x)}{dx} + m x P_l^m(x) = -\sqrt{1 - x^2}P_l^{m+1}(x), \\
		&(1 - x^2)\frac{dP_l^m(x)}{dx} - m x P_l^m(x) = (l+m)(l-m+1)\sqrt{1 - x^2}P_l^{m-1}(x), \\
		&(1 - x^2)\frac{dP_l^m(x)}{dx} = \frac{1}{2l+1}\left(-l(l-m+1)P_{l+1}^m(x)+(l+1)(l+m)P_{l-1}^m(x)\right).
	\end{align*}
	
	根据球谐函数与连带勒让德函数的关系：
	\[
	Y_l^m(\theta,\phi) = \sqrt{\frac{2l+1}{4\pi}\frac{(l-m)!}{(l+m)!}} P_l^m(\cos\theta)e^{im\phi}.
	\]
	
	综合以上积分式，进行$x=\cos\theta$换元，易得：
	\begin{align*}
		&\int_0^\pi\int_0^{2\pi} \overline{Y_{L}^{M}(\theta, \phi)}e^{i\phi}\left(\frac{\partial Y_{l}^{m}(\theta,\phi)}{\partial\theta}\sin\theta\cos\theta+i\frac{\partial Y_{l}^{m}(\theta,\phi)}{\partial\phi}\right) d\phi d\theta\\
		&=-\frac{1}{2}\delta_{M,m+1}\sqrt{\frac{(2l+1)(2L+1)(L-m-1)!(l-m)!}{(L+m+1)!(l+m)!}}\int_{-1}^{1}\left(x\sqrt{1-x^2}\frac{dP_l^m(x)}{dx}+\frac{m}{\sqrt{1-x^2}}P_l^m(x)\right)P_L^{m+1}(x)\,dx\\[6pt]
		&=\delta_{M,m+1}\left(\delta_{L,l+1}\,l\sqrt{\frac{(l+m+1)(l+m+2)}{(2l+1)(2l+3)}}+\delta_{L,l-1}\,(l+1)\sqrt{\frac{(l-m)(l-m-1)}{4l^2-1}}\right),\\[12pt]
		%
		&\int_0^\pi\int_0^{2\pi}  \overline{Y_{L}^{M}(\theta, \phi)}e^{-i\phi}\left(\frac{\partial Y_{l}^{m}(\theta,\phi)}{\partial\theta}\sin\theta\cos\theta - i\frac{\partial Y_{l}^{m}(\theta,\phi)}{\partial\phi}\right) d\phi d\theta\\
		&=-\frac{1}{2}\delta_{M,m-1}\sqrt{\frac{(2l+1)(2L+1)(L-m+1)!(l-m)!}{(L+m-1)!(l+m)!}}\int_{-1}^{1}\left(x\sqrt{1-x^2}\frac{dP_l^m(x)}{dx}-\frac{m}{\sqrt{1-x^2}}P_l^m(x)\right)P_L^{m-1}(x)\,dx\\[6pt]
		&=-\delta_{M,m-1}\left(\delta_{L,l+1}\,l\sqrt{\frac{(l-m+1)(l-m+2)}{(2l+1)(2l+3)}}+\delta_{L,l-1}\,(l+1)\sqrt{\frac{(l+m)(l+m-1)}{4l^2-1}}\right),\\[12pt]
		%
		&\int_0^\pi\int_0^{2\pi}  \overline{Y_{L}^{M}(\theta, \phi)}\frac{\partial Y_{l}^{m}(\theta,\phi)}{\partial\theta}\sin^2\theta\, d\phi d\theta\\
		&=-\frac{1}{2}\delta_{M,m}\sqrt{\frac{(2l+1)(2L+1)(L-m)!(l-m)!}{(L+m)!(l+m)!}}\int_{-1}^{1}(1-x^2)\frac{dP_l^m(x)}{dx}P_L^m(x)\,dx\\[6pt]
		&=\delta_{M,m}\left(\delta_{L,l+1}\,l\sqrt{\frac{(l+1)^2-m^2}{(2l+1)(2l+3)}}-\delta_{L,l-1}\,(l+1)\sqrt{\frac{l^2-m^2}{4l^2-1}}\right).
	\end{align*}
	
	即证题式。
	
	另外，球坐标系梯度算符在直角坐标系下的分量为：
	\begin{align*}
		r\nabla &= \frac{\hat{\bm{e}}_x - i\hat{\bm{e}}_y}{2}e^{i\phi}\left(\cos\theta\frac{\partial}{\partial\theta}+\frac{i}{\sin\theta}\frac{\partial}{\partial\phi}\right)+\frac{\hat{\bm{e}}_x + i\hat{\bm{e}}_y}{2}e^{-i\phi}\left(\cos\theta\frac{\partial}{\partial\theta}-\frac{i}{\sin\theta}\frac{\partial}{\partial\phi}\right)-\hat{\bm{e}}_z\sin\theta\frac{\partial}{\partial\theta},
	\end{align*}
	
	根据如上连带勒让德函数及其导数的递推式，易知，其作用于球谐函数 $Y_l^m(\theta,\phi)$ 得：
	\begin{align*}
		r\nabla Y_l^m(\theta,\phi)&=\left(-l\sqrt{\frac{(l+1)^2-m^2}{(2l+1)(2l+3)}}Y_{l+1}^m+(l+1)\sqrt{\frac{l^2-m^2}{4l^2-1}}Y_{l-1}^m\right)\hat{\bm{e}}_z\\[6pt]
		&\quad+\left(l\sqrt{\frac{(l+m+1)(l+m+2)}{(2l+1)(2l+3)}}Y_{l+1}^{m+1}+(l+1)\sqrt{\frac{(l-m)(l-m-1)}{4l^2-1}}Y_{l-1}^{m+1}\right)\frac{\hat{\bm{e}}_x - i\hat{\bm{e}}_y}{2}\\[6pt]
		&\quad-\left(l\sqrt{\frac{(l-m+1)(l-m+2)}{(2l+1)(2l+3)}}Y_{l+1}^{m-1}+(l+1)\sqrt{\frac{(l+m)(l+m-1)}{4l^2-1}}Y_{l-1}^{m-1}\right)\frac{\hat{\bm{e}}_x + i\hat{\bm{e}}_y}{2}.
	\end{align*}
	
	下根据球谐函数正交归一性即可，略。
\end{document}
